\documentclass[10pt]{extarticle}

% ---------- Layout ----------
\usepackage[a4paper,margin=1.35cm,columnsep=0.6cm,top=1.15cm,bottom=1.15cm]{geometry}
\usepackage{multicol}
\twocolumn
\usepackage[T1]{fontenc}
\usepackage[utf8]{inputenc}
\usepackage{lmodern}
\usepackage{microtype}
\usepackage{enumitem}
\setlist{topsep=1pt,itemsep=1pt,parsep=0pt,partopsep=0pt,leftmargin=*}
\usepackage{titlesec}
\usepackage{xcolor}
\definecolor{hdrblue}{RGB}{20,55,95}
\definecolor{rule}{RGB}{20,55,95}
\titleformat{\section}{\normalfont\sffamily\bfseries\color{hdrblue}\large}{\thesection.}{0.4em}{}
\titlespacing*{\section}{0pt}{5pt}{2pt}
\titleformat{\subsection}{\normalfont\sffamily\bfseries\small}{\thesubsection}{0.4em}{}
\titlespacing*{\subsection}{0pt}{3pt}{1pt}
\setlength{\parindent}{0pt}
\setlength{\parskip}{2pt}
\usepackage[colorlinks=true,linkcolor=hdrblue,urlcolor=hdrblue,citecolor=hdrblue,
            pdftitle={CS455 Deliverable 0 -- Project Proposal},
            pdfauthor={Prithviraj Ghosh, Mainak Sarkar, Khushi Ranawat, Mahir Jain, Chaitanya Nitawe}]{hyperref}
\usepackage{url}
\usepackage{fancyhdr}
\pagestyle{fancy}
\fancyhf{}
\renewcommand{\headrulewidth}{0pt}

\begin{document}

% ---------- Title block ----------
\twocolumn[
  \begingroup
  \centering
  {\sffamily\bfseries\LARGE\color{hdrblue}
    Design and Implementation of a Campus-Centric Carpooling System}\par
  \vspace{3pt}

  {\small
    Prithviraj Ghosh (230794) \quad
    Mainak Sarkar (230619) \quad
    Khushi Ranawat (230562) \quad
    Mahir Jain (230618) \quad
    Chaitanya Nitawe (230315)}\par
  \endgroup
  \vspace{2pt}
  {\color{rule}\hrule height 0.8pt}
  \vspace{4pt}
]

% ============================================================
\section{Project Title \& Description}
\textbf{Campus Ride-Pooling and Split-Fare System.}
Students travelling from campus to nearby transit hubs (railway station,
airport, etc.) currently book cabs individually and pay full fare, even
when several people are headed to the same place at the same time. This
project is a web platform that matches such riders on overlapping
departure windows, allocates them seats in a shared vehicle, splits the
fare, and coordinates pickup, with a governed AI agent assisting in the
workflow.

% ============================================================
\section{User Roles}

\textbf{Rider / Ride Owner:} Creates, modifies, or cancels rides;
requests to join existing rides; accepts/rejects co-rider join requests
when acting as ride owner; pays their fare share; rates co-riders.

\textbf{Operations Admin:} Manages registered vehicles and vehicle
capacity; configures destination hubs and fare bands; handles SOS
incidents; reviews the AI agent's action log; issues warnings or
suspends accounts.

\textbf{Authentication:} IITK email-based authentication for riders,
with separate role-based authorization for Operations Admin access.

% ============================================================
\section{Core Transactional Workflow}
Following IITK email authentication, new users complete a quick setup,
selecting from a pre-fixed set of tags/interests (e.g., academic
tracks, hobbies, or ``quiet ride''); this data directly feeds the AI
matchmaking agent.

A user initiates a ride with a departure time and preferred vehicle
(cab or auto), or requests to join an existing one from AI-generated
suggestions based on similar start/destination
$\rightarrow$ co-riders join pending ride-owner approval
$\rightarrow$ the fare split is decided (re-computed if a new rider
joins dynamically) $\rightarrow$ the trip executes and completes.

Once the dynamic fare split locks in, a temporary real-time in-app
chat automatically opens for the active pool to coordinate the exact
physical meetup spot at the campus gates. Post-ride, users settle the
fare and may file detailed text complaints or rate co-riders; these
reports are routed immediately to the AI Complaint Management system
for severity scoring and potential disciplinary action.

% ============================================================
\section{Shared, Contended Resource \& Concurrency Challenge}

\textbf{(i) Vehicle Capacity -- Stale Availability:} The core
concurrency challenge is ensuring riders do not join a ride based on
outdated seat availability. E.g., a rider may see one seat available
and attempt to join while the ride owner concurrently accepts a
different pending rider for that same seat; without proper
concurrency control the first request could still be accepted on
stale data, causing overbooking. The system must therefore
\emph{revalidate seat availability and update the ride atomically}
when a rider is accepted.

\textbf{(ii) Time of Departure -- Dynamic Itinerary Drift:} Adding a
new rider introduces pickup detours that shift departure and arrival
times for everyone. If multiple users apply concurrently against an
outdated departure time, the global ETA can drift mid-transaction,
causing missed pickup windows and itinerary conflicts.

% ============================================================
\section{AI Agent}

\textbf{(i) Pool Car Recommendation System:} An AI agent, activated
via an ``Ask AI to Suggest'' button, ranks pool cars for searching
riders, generating personalized pros/cons based on pricing, route
efficiency, shared-profile interests, user reviews, and past
complaints. For incoming join requests, the agent gives the ride
owner a compatibility summary based on relevant ride preferences and
logistics. An interactive AI chatbot additionally lets users filter
rides, inspect co-rider details, and manage pool preferences via
natural-language queries.

\textbf{(ii) Complaint Management System:} The AI summarizes
complaints, assigns a severity score, and categorizes them with
behavioral tags (e.g., safety, tardiness). By recognizing recurring
patterns in a user's history, it proactively generates enforcement
triggers -- warnings or suggested suspensions -- routed to an Admin
Dashboard for quick, one-click human approval.

% ============================================================
\section{Proposed Technology Stack}
\begin{itemize}
  \item TypeScript
  \item Node.js (backend)
  \item Firestore / PostgreSQL
  \item LLM API (free tier)
  \item React.js / Next.js
  \item Socket.IO (real-time communication)
  \item GCP -- Cloud Run, Cloud SQL, Cloud Scheduler, Cloud Logging
\end{itemize}

% ============================================================
\section{Initial Risks \& Technical Uncertainties}
\begin{itemize}
  \item Concurrency correctness under simultaneous seat and ownership actions
  \item LLM hallucination in recommendation and complaint scoring
  \item API and cloud infrastructure cost management
  \item Designing a robust and fair complaint-management pipeline
\end{itemize}

% ============================================================
\section{Jira / GitHub}
\textbf{GitHub repository:} \\
{\tiny\url{https://github.com/mainak9093/Campus_Ridepooling_Automation}}
\ -- Ride-pooling and split-fare platform for campus travel: matches
riders on overlapping departure windows, allocates shared vehicle
seats under concurrency control, and splits fares.

\textbf{Jira project:} \\
{\tiny\url{https://cs455project.atlassian.net/jira/software/projects/SCRUM/boards/1}}

Instructor and TAs will be added to both the GitHub repository and
the Jira project prior to submission.

\end{document}
